% GENERATED by scripts/build.py convert ch03 -- do not hand-edit; edit
% chapters/md/03_data_and_variables.md and regenerate.
% md sha256 aba6024df81e339869a9d46ff5a4287955bf29eed058b70cc55260fa10dee3a6
% canonical manuscript/v13_2026-10-17/03_data_and_variables.md sha256 aba6024df81e339869a9d46ff5a4287955bf29eed058b70cc55260fa10dee3a6
% md identical to canonical: True
% filter sha256 d487b0e18e456f08097065d7ceaaf39fef0daee29958941b5539c501e0e4c24a

\section{Data and variable construction}\label{sec:3}\ThesisCheckNumber{section}{3}

This section documents the baseline and historical samples separately: they come from different DHS rounds, use different exposure constructions, and are validated independently, so no claim about one sample applies to the other. Full variable definitions and sample-flow tables are in Appendix A; estimation results are reported in Section 5.

\subsection{Baseline sample: unit of analysis, rounds, and pooling}\label{sec:3.1}\ThesisCheckNumber{subsection}{3.1}

The unit of analysis is the child. The baseline pools each country's most recent DHS round as of data collection: Ethiopia (2024--25), Ghana (2022), Kenya (2022), Nigeria (2024). The pooled, harmonised child file contains 70,231 children across 4,481 DHS clusters and 114 administrative-region fixed-effect cells (one region set per country). This pooled file is \textbf{not} the regression sample for any model; it is the universe before outcome-validity and exposure-classifiability restrictions are applied.

\begin{figure}[H]
\centering
\includegraphics[width=0.92\textwidth]{\ThesisRoot/figures/map_survey_cluster_coverage.png}
\ThesisFigureCaption{3.1}{Survey cluster locations, by country (descriptive only; coordinates randomly displaced for confidentiality, Section 3.4).}
\end{figure}

\subsection{Baseline anthropometric outcomes and under-five eligibility}\label{sec:3.2}\ThesisCheckNumber{subsection}{3.2}

Three outcomes follow standard DHS/WHO practice: height-for-age (HAZ), weight-for-age (WAZ), weight-for-height (WHZ), each a z-score against the WHO Child Growth Standards \autocite{deOnis2006}. A child contributes to an outcome only if its z-score is flagged valid in the DHS recode (excluding DHS's own special codes 9996--9999 and system-missing values). Valid measurements give smaller outcome-specific samples: 36,985 (HAZ), 37,260 (WAZ), 37,088 (WHZ). Requiring both household and community exposures to be classifiable for the same child gives common samples of 36,040 (HAZ), 36,308 (WAZ), 36,145 (WHZ), across 4,426 clusters. Age is taken uniformly from DHS's standard \ThesisCode{b19} field for all four countries; the historical extension (Sections 3.5--3.8) builds its own, differently constructed \ThesisCode{age\_\allowbreak{}m} field and the two must not be conflated.

\begin{figure}[H]
\centering
\includegraphics[width=0.92\textwidth]{\ThesisRoot/figures/figure1_prevalence_by_country.png}
\ThesisFigureCaption{3.2}{Weighted anthropometric prevalence by country, with 95\% confidence intervals (each outcome's own outcome-valid sample, Section 3.2).}
\end{figure}

Figure 3.2 reports weighted stunting, wasting and underweight prevalence by country, each estimated on that outcome's own outcome-valid sample above, under a stratum/PSU/weight survey design (\ThesisCode{strata}, \ThesisCode{psu}, the \ThesisCode{v005}-derived weight of Section 3.4) with 95\% confidence intervals from a logit-transformed interval. This descriptive design is distinct from the PSU-clustered inference the regression models in Section 4 implement; the two are not interchangeable. Stunting is highest in Ethiopia (39.9\%) and Nigeria (39.1\%), roughly double Ghana (16.9\%) and Kenya (17.4\%); wasting is lower and more similar across countries (4.9--8.4\%); underweight is highest in Nigeria (26.4\%) and lowest in Kenya (10.0\%). These are descriptive estimates, not regression results; Section 5 models the within-country variation they motivate.

\subsection{Baseline WASH exposures: household and leave-one-out community}\label{sec:3.3}\ThesisCheckNumber{subsection}{3.3}

The baseline sample has four exposures: own-household improved/non-improved water and sanitation, and their leave-one-out community counterparts, following the WHO/UNICEF JMP ``improved'' classification (the same classification independently verified for the historical extension's gridded product, Section 3.5, Appendix A). Household WASH is a household-level binary indicator. \textbf{Community WASH is an unweighted leave-one-out proportion among other sampled, classifiable HR households in the same DHS cluster, excluding the child's own household} --- not a census of the community, and not an average over other children. Households recorded as Unknown or not applicable, and (Ghana's core sanitation variant) an unclassified biodigester, are excluded from both numerator and denominator. The resulting rate is a \([0,1]\) proportion, missing if no other classifiable household exists in the cluster; coefficients on it are therefore per full unit of coverage (0 to 1), not per 10 percentage points as in the historical tables of Section 5.6. Improved shared sanitation counts as improved via a Basic+Limited collapse; a robustness variant treats Ghana's biodigester differently (Appendix A, Table A.1). A child's own classification is never part of its own community exposure, which makes community WASH vary mainly \emph{between} clusters --- Section 4.3 derives why this rules out a cluster-fixed-effects specification for this exposure, unlike for household WASH.

\begin{figure}[H]
\centering
\includegraphics[width=0.78\textwidth]{\ThesisRoot/figures/figure3_wash_characteristics.png}
\ThesisFigureCaption{3.3}{Weighted household WASH characteristics by country, as percentages of each country's full weighted sample (Unknown and structurally-not-applicable cases included in the denominator, not excluded). The sanitation panel uses the collapsed Improved/Non-improved scale: open defecation and other unimproved facility types are folded into Non-improved, not shown as a separate segment.}
\end{figure}

Figure 3.3 situates these exposures in their country context. Its denominator is each country's full weighted descriptive sample, not the classifiable-household denominator used to construct the regression's leave-one-out community exposure (Section 3.3) --- a different quantity for a different purpose. The left panel shows drinking water at full category detail, with surface water and unimproved as their own segments; the right panel shows sanitation on the collapsed Improved/Non-improved scale used in the regressions, where open defecation and other unimproved types are folded into Non-improved rather than displayed separately. Unknown, structurally-not-applicable and (sanitation only) unclassified-biodigester cases remain their own segments in both panels. Improved water ranges from 68.1\% (Ethiopia) to 84.2\% (Ghana). Improved sanitation is lower and ordered differently: 42.5\% (Ethiopia), 55.2\% (Ghana), 61.2\% (Nigeria), 63.8\% (Kenya). Only Ethiopia combines high stunting (Figure 3.2) with low sanitation coverage; Nigeria has the second-highest stunting prevalence but sanitation coverage close to Kenya's, so the two figures together do not support a simple cross-country ordering linking WASH coverage to stunting. Sections 4 and 5 examine the within-country, household- and community-level association directly, which this descriptive comparison cannot.

\subsection{Baseline covariates, geography, and weights}\label{sec:3.4}\ThesisCheckNumber{subsection}{3.4}

Covariates include child sex, age in months (\ThesisCode{b19}), maternal age, maternal education, literacy, household size, birth order, wealth quintile, and urban/rural residence. The baseline is built from country-specific DHS recode files (the Children's Recode, KR), not IPUMS-DHS, using DHS's own field names throughout --- distinct from the historical extension's IPUMS-harmonised variables (Sections 3.5--3.8). The KR file's weight field is \ThesisCode{v005}, the woman's individual sample weight (rescaled by the documented 1,000,000 factor), not the household weight \ThesisCode{hv005}; this is not \ThesisCode{PERWEIGHT}, the different field the historical extension uses. The cluster identifier \ThesisCode{v021}, combined with a country prefix into \ThesisCode{country\_\allowbreak{}psu}, is analogous in purpose to the historical extension's \ThesisCode{IDHSPSU} but constructed independently. DHS cluster coordinates are randomly displaced for confidentiality (urban up to 2 km; rural 99\% up to 5 km, 1\% up to 10 km) \autocite{Burgert2013} --- inconsequential for the baseline models, which use no coordinates directly, but consequential for the historical extension's spatial matching (Section 3.5). Baseline models use the \ThesisCode{v005}-derived weight, normalised so each country contributes equal total weight mass regardless of raw sample size (Section 4.5 gives the formula).

\subsection{Historical extension: rounds, exposure source, and scope}\label{sec:3.5}\ThesisCheckNumber{subsection}{3.5}

The historical extension links an earlier DHS round for three countries --- Ghana (2014), Kenya (2014), Nigeria (2018) --- to IHME's annual gridded estimates of local access to improved water (W\_IMP) and sanitation (S\_IMP), available 2000--2017 \autocite{LBDWaSH2020}, verified directly from the primary source (Appendix A, Table A.2). The extraction computes an \textbf{area average of the raster surface within each child's exposure buffer}, not the source study's population-weighted aggregate --- a distinction restated wherever a historical coefficient is reported. Ethiopia's 2016 round enters only as a separately labelled, provisional sensitivity analysis (Appendix C), because its Ethiopian-calendar-to-Gregorian interview-date conversion could not be validated to the same standard as the three established-date rounds. The interview date does not place a child's exposure window --- windows are placed from the birth date alone (Section 3.7; Appendix C) --- but governs eligibility: date validity, birth preceding interview, and (for the two postnatal windows) whether the window had completed by interview. This is a measurement improvement on timing specifically, not a record of a child's actual household facilities, which no available data source can provide.

\subsection{Historical spatial extraction: buffers, displacement, and coverage}\label{sec:3.6}\ThesisCheckNumber{subsection}{3.6}

Each child's exposure buffer is centred on its cluster's displaced coordinates (azimuthal equidistant projection). The primary buffer radius is 2 km (urban) or 5 km (rural). Displacement describes locational \emph{uncertainty}; the buffer radius defines the \emph{area averaged}. A labelled sensitivity repeats the rural assignment at 10 km, and a separate sensitivity uses containing-pixel assignment instead (Section 5.8, Appendix C.4). The raster value is averaged by closed-form disk--polygon overlap, not centroid or nearest-neighbour approximation. A child-year is assigned a valid extraction only when its valid-pixel area meets a documented tolerance (within \(10^{-3}\) of full coverage); partial- or missing-coverage child-years are excluded rather than assigned from incomplete coverage.

\subsection{Historical exposure windows and completion rules}\label{sec:3.7}\ThesisCheckNumber{subsection}{3.7}

Four candidate windows apply to the same annual raster series. \textbf{The primary window is the first 12 completed postnatal months} (birth month through birth month + 11, with completion requiring the interview month \(t \geq\) birth month + 12; Section 6.4 states why this window was proposed). \textbf{The first 24 postnatal months} (birth month through birth month + 23, completion \(t \geq\) birth month + 24) is a sensitivity window, not the primary one. Two further alternatives carry no postnatal completion rule: the prenatal window (the nine months before birth, ending at birth) and the birth-year window (the full calendar year of birth). A window spanning two calendar years is weighted by each year's share of months, and every touched year must have a verified raster value. The birth-year window's lack of a completion check is a disclosed limit on that window's eligibility rule, not an oversight. Because IHME's product is annual, the value used for a window's final touched year can reflect conditions after the interview even when the postnatal completion check is satisfied; this project flags every such case but does not exclude or adjust for it.

\subsection{Historical covariates, weights, and the established-date measurement universe}\label{sec:3.8}\ThesisCheckNumber{subsection}{3.8}

Historical models control for cluster fixed effects, survey-qualified country-by-birth-year terms, child sex, maternal age, maternal education, marital status, and completed age in months. Birth-year terms are each country's own birth-year effects relative to its own earliest observed year, not one pooled reference year; Section 4.6 shows this choice does not itself change the fitted exposure coefficient. The age control uses its own \ThesisCode{age\_\allowbreak{}m} field (reported months where the source delivers them, Ethiopia/Nigeria; a calendar-month-difference construction otherwise, Ghana/Kenya), distinct from the baseline's \ThesisCode{b19}. Wealth quintile is \textbf{excluded from the main specification}, because current wealth may be a post-exposure mediator of historical WASH conditions while also being correlated with the exposure itself --- a genuine trade-off, not a one-sided argument: current wealth could instead proxy \emph{prior} socioeconomic circumstances, in which case excluding it leaves that confounding unaddressed. Wealth is included only as a labelled sensitivity (Section 5.8), where the pooled water--HAZ coefficient is essentially unchanged (1.01 vs.~1.03). This differs from, and is more defensible than, claiming cluster fixed effects absorb wealth mechanically: they absorb only what is constant \emph{within} a cluster, confirmed directly since the baseline cluster-fixed-effects model includes and estimates wealth as a standard covariate (Section 4.2).

Weighting uses PERWEIGHT, normalised so each of the three established-date countries contributes an equal total weight within each model's own sample, reported alongside its unweighted counterpart for the headline cell (Section 5.6). No DHS or IPUMS source documents a sampling weight specific to the measured-anthropometry subsample (Appendix A); the working inference, stated wherever historical coefficients are reported, is that within-survey PERWEIGHT normalisation removes a constant measurement-subsample selection factor under random household-level selection --- it cannot correct for differential non-response among eligible but unmeasured children. Historical results are therefore reported as \textbf{weighted conditional associations among eligible measured children}, not population-representative estimates.

\subsection{A measurement-universe discrepancy: Nigeria 2018}\label{sec:3.9}\ThesisCheckNumber{subsection}{3.9}

Three distinct child populations arise for Nigeria 2018. The official 2018 NDHS report counts 12,806 (unweighted) children under five eligible and measured (men's-questionnaire subsample households). The IPUMS-DHS birth-history extract used here contains 33,924 birth-history records, of which 11,704 are non-missing for anthropometry. The final historical estimation sample (water, HAZ, post-12-month window) retains 6,912 after window-completion, spatial-coverage and singleton-cluster exclusions. These differ by construction --- official eligible count, interviewed-women's birth-history children, and the further-restricted estimation sample are three different populations --- and the gap between the official count and the extract (12,806 vs.~11,704) is disclosed as an open, unresolved question, not treated as resolved by any estimation choice here.
